\section{課題}

「C言語（中級）」で与えられている {\tt hash.c} を参考に，次の仕様を満たすハッシュ表を C++ の class で実現する．

{\tt ()}，{\tt insert}，{\tt -=} の 3 つのメンバ関数で，キー（文字列）による検索，アイテムの挿入，削除を行える表を C++ で実現し，振舞いを確認する．ハッシュ表オブジェクトを {\tt table} とすると，検索は {\tt table("hoge")}，挿入は {\tt table.insert("hoge", x)}，削除は {\tt table -= "hoge"} のように操作するものとする．

また，加点課題として，いろいろな型の値を収納できるよう（キーは文字列でよい），クラスをテンプレートにする．
